% fig_vl20_mikro.tex — Die Energieschale: das Abzählen der Mikrozustände
% in [E, E+\Delta], und was aus dem Zählen wird, wenn N wächst.
%
% WHY THIS FIGURE.  The lecture writes down
%       rho = sum_{E_j in [E,E+Delta]} (1/N_E) |E_j><E_j|,   S = k log(N_E)
% and makes N_E concrete only on a three-level toy, where N_E = 2.  Two
% things stay invisible there, and both are shapes rather than sentences:
%   (i)  what "counting the states in the shell" looks like when the
%        states are degenerate — the count is a binomial coefficient, and
%        the multiplicities are wildly unequal from level to level;
%   (ii) that this toy count IS the thermodynamic entropy: k log(N_E)
%        divided by N converges to a smooth concave curve of the energy
%        per particle, with maximum log 2, whose slope is what the next
%        lecture calls 1/T.
% Panel (a) can be counted with the finger; panel (b) says that the same
% count, done for 10^23 spins, is the entropy of thermodynamics.
%
% GEOMETRY FROM THE FORMULA.  Nothing is drawn by eye.  The model is N
% non-interacting spins in a field, single-spin energies 0 and epsilon,
%
%       H |s> = k epsilon |s>,   k = number of up spins in s,
%
% so the spectrum is the equally spaced ladder E_k = k epsilon, and the
% degeneracy of level k is the number of strings with k up spins,
%
%       g_k = C(N,k).
%
% Panel (a), N = 6:  g = (1,6,15,20,15,6,1), sum = 2^6 = 64 dots, one dot
% per microstate, placed at height E_k = k epsilon and spread horizontally
% at unit spacing about x = 0 — so the width of a row IS its degeneracy
% and the rows may be compared by eye.  The shell is E = 1.5 epsilon,
% Delta = 2 epsilon, i.e. [1.5, 3.5] epsilon, which admits exactly the
% levels k = 2 and k = 3:
%
%       N_E = C(6,2) + C(6,3) = 15 + 20 = 35,   p_j = 1/35,
%       S   = k log(35).
%
% Deliberately two levels, not one: N_E counts eigenvalues in an interval,
% and the interval need not respect the level structure — that is exactly
% what the lecture's own example does (E_2 and E_3 both fall in [0.9,1.2]).
%
% Panel (b), the same count for a shell holding one level, per spin:
%
%       S/(kN) = (1/N) log C(N,k),   e = E/(N epsilon) = k/N,
%
% plotted from exact binomial coefficients for N = 6 and N = 24, against
% its N -> infinity limit, the binary entropy
%
%       s(e) = -e log(e) - (1-e) log(1-e),
%
% whose maximum s(1/2) = log 2 is the ceiling S <= k log(d) of the
% previous lecture, here d = 2^N.  The finite-N clouds sit below the
% curve and creep up towards it: the correction is -log(N)/(2N), which is
% why N = 6 is already the right shape and 10^23 is indistinguishable
% from the limit.
%
% DESIGN DECISIONS
%   * colour is probability, and nothing else: cgblue = the states that
%     actually occur (p_j = 1/N_E), grey open circles = the states with
%     p_j = 0.  The same filled/open distinction is repeated in (b), where
%     the two filled points are precisely the two levels of the shell of
%     (a) — that is the link between the panels, and it needs no words;
%   * energies outside the shell are hatched, not filled, exactly as the
%     classically forbidden region of fig_vl11 is: unreachable is
%     unreachable, and hatching survives greyscale;
%   * the energy levels themselves are neutral grey dashed rules; the two
%     shell edges are solid spacecadet, because they are the one piece of
%     input information the ensemble is built from;
%   * in (b) the finite sizes are dot clouds (the energy really is
%     discrete there, e = k/N) and only the thermodynamic limit is drawn
%     as a continuous curve.  The three sizes N = 6, 24, infinity are the
%     ordered ramp cgblue -> birrenorange -> birrenred of `tpclevels`;
%     munsell is skipped for N = 24 for the reason fig_vl11 gives — beside
%     cgblue it is the same teal, so it could not carry the code.  N = 6
%     keeps the cgblue of panel (a), because it IS the system of panel (a);
%   * (b) plots shells that hold a single level, so there N_E = C(N,k),
%     whereas the shell of (a) straddles two.  The corner of (b) says so:
%     otherwise the reader would expect log(35) where log(20) is drawn;
%   * the degeneracy column on the right of (a) is ruled off and summed,
%     so that sum_k g_k = 2^N = dim(H) is on the plate: the shell keeps 35
%     of 64, and the reader can check both numbers by counting dots;
%   * every curve is labelled where it is, no legend box.

\begin{center}
\begin{tikzpicture}

% ─────────────────────────────────────────────────────────────
%  (a)  All 2^6 microstates, sorted by energy; the shell picks two rows
% ─────────────────────────────────────────────────────────────
\begin{axis}[
    name=schale,
    tpclevelplot,
    width=7.4cm, height=4.8cm, scale only axis,
    xmin=-15.4, xmax=13.8,
    ymin=-2.5,  ymax=7.3,
    title={(a) \tpcifenglish{energy shell of $N=6$ spins}%
                             {Energieschale von $N=6$ Spins}},
]

  % Energies outside [E, E+Delta]: p_j = 0, hatched like a forbidden region
  \fill[forbidden] (axis cs:-11,-0.55) rectangle (axis cs:11,1.5);
  \fill[forbidden] (axis cs:-11,3.5)   rectangle (axis cs:11,6.55);

  % The ladder E_k = k epsilon: neutral reference rules
  \draw[levellinegrey] (axis cs:-10.6,0) -- (axis cs:10.6,0);
  \draw[levellinegrey] (axis cs:-10.6,1) -- (axis cs:10.6,1);
  \draw[levellinegrey] (axis cs:-10.6,2) -- (axis cs:10.6,2);
  \draw[levellinegrey] (axis cs:-10.6,3) -- (axis cs:10.6,3);
  \draw[levellinegrey] (axis cs:-10.6,4) -- (axis cs:10.6,4);
  \draw[levellinegrey] (axis cs:-10.6,5) -- (axis cs:10.6,5);
  \draw[levellinegrey] (axis cs:-10.6,6) -- (axis cs:10.6,6);

  % Which way is up: energy increases along the page
  \draw[axisstyle] (axis cs:-14.4,4.25) -- (axis cs:-14.4,6.5);
  \node[lbl, text=birrengray, font=\scriptsize, rotate=90]
    at (axis cs:-13.5,5.35) {\tpcifenglish{energy}{Energie}};

  % The two edges of the shell — the whole input to the ensemble
  \draw[potential] (axis cs:-11,1.5) -- (axis cs:11,1.5);
  \draw[potential] (axis cs:-11,3.5) -- (axis cs:11,3.5);
  \node[lbl, anchor=east] at (axis cs:-11.3,1.5) {$E$};
  \node[lbl, anchor=east] at (axis cs:-11.3,3.5) {$E+\Delta$};

  % Microstates outside the shell: p_j = 0
  \addplot[only marks, mark=o, mark size=1.25pt,
           mark options={draw=birrengray, fill=white, line width=0.4pt}]
    coordinates {
      (0.0,0) (-2.5,1) (-1.5,1) (-0.5,1) (0.5,1) (1.5,1) (2.5,1)
      (-7.0,4) (-6.0,4) (-5.0,4) (-4.0,4) (-3.0,4) (-2.0,4) (-1.0,4)
      (0.0,4) (1.0,4) (2.0,4) (3.0,4) (4.0,4) (5.0,4) (6.0,4) (7.0,4)
      (-2.5,5) (-1.5,5) (-0.5,5) (0.5,5) (1.5,5) (2.5,5) (0.0,6)
    };

  % Microstates inside the shell: p_j = 1/N_E
  \addplot[only marks, mark=*, mark size=1.35pt,
           mark options={draw=cgblue, fill=cgblue}]
    coordinates {
      (-7.0,2) (-6.0,2) (-5.0,2) (-4.0,2) (-3.0,2) (-2.0,2) (-1.0,2)
      (0.0,2) (1.0,2) (2.0,2) (3.0,2) (4.0,2) (5.0,2) (6.0,2) (7.0,2)
      (-9.5,3) (-8.5,3) (-7.5,3) (-6.5,3) (-5.5,3) (-4.5,3) (-3.5,3)
      (-2.5,3) (-1.5,3) (-0.5,3) (0.5,3) (1.5,3) (2.5,3) (3.5,3)
      (4.5,3) (5.5,3) (6.5,3) (7.5,3) (8.5,3) (9.5,3)
    };

  % Probabilities, written where they hold
  \node[lbl, text=cgblue] at (axis cs:0,2.5) {$p_j=1/35$};
  \node[lbl, text=birrengray, fill=white, inner sep=1.5pt]
    at (axis cs:0,5.5) {$p_j=0$};
  \node[lbl, text=birrengray, fill=white, inner sep=1.5pt]
    at (axis cs:0,0.5) {$p_j=0$};

  % Degeneracies g_k = C(6,k), read off as the width of each row
  \node[lbl, text=birrengray, font=\scriptsize, anchor=west]
    at (axis cs:11.25,6.9) {$g_k=\tbinom{6}{k}$};
  \node[lbl, text=birrengray, anchor=west] at (axis cs:11.3,0) {$1$};
  \node[lbl, text=birrengray, anchor=west] at (axis cs:11.3,1) {$6$};
  \node[lbl, text=cgblue,     anchor=west] at (axis cs:11.3,2) {$15$};
  \node[lbl, text=cgblue,     anchor=west] at (axis cs:11.3,3) {$20$};
  \node[lbl, text=birrengray, anchor=west] at (axis cs:11.3,4) {$15$};
  \node[lbl, text=birrengray, anchor=west] at (axis cs:11.3,5) {$6$};
  \node[lbl, text=birrengray, anchor=west] at (axis cs:11.3,6) {$1$};

  % The column of degeneracies adds up to the dimension of the Hilbert space
  \draw[birrengray!70, line width=0.4pt]
    (axis cs:11.2,-0.5) -- (axis cs:12.9,-0.5);
  \node[lbl, text=birrengray, font=\scriptsize, anchor=north west]
    at (axis cs:11.2,-0.55) {$2^6=64$};

  % What the model is, and what one dot means
  \node[lbl, text=birrengray, font=\scriptsize, anchor=west]
    at (axis cs:-15.4,7.0)
    {$E_k=k\varepsilon$; \tpcifenglish{one dot $=$ one microstate}%
                                     {ein Punkt $=$ ein Mikrozustand}};

  % The count itself
  \node[lbl, anchor=north] at (axis cs:-0.7,-0.95)
    {$N_E=\tbinom{6}{2}+\tbinom{6}{3}=15+20=35$};
  \node[lbl, anchor=north] at (axis cs:-0.7,-1.80)
    {$S=k\log(N_E)=k\log(35)$};
\end{axis}

% ─────────────────────────────────────────────────────────────
%  (b)  The same count per spin, and its thermodynamic limit
% ─────────────────────────────────────────────────────────────
\begin{axis}[
    at={(schale.south east)}, anchor=south west, xshift=20mm,
    tpcplot,
    width=4.2cm, height=4.8cm, scale only axis,
    xmin=-0.03, xmax=1.03, ymin=0, ymax=0.79,
    xtick={0,0.5,1}, xticklabels={$0$,$\tfrac12$,$1$},
    ytick={0,0.2,0.4,0.6},
    xlabel={$e=E/(N\varepsilon)$},
    ylabel={$S/(kN)$},
    every axis x label/.style={at={(0.5,-0.115)}, anchor=north, font=\small},
    every axis y label/.style={at={(axis description cs:-0.17,0.5)},
                               rotate=90, anchor=south, font=\small},
    title={(b) \tpcifenglish{entropy per spin}{Entropie pro Spin}},
]
  % Ceiling S <= k log(d), d = 2^N: one log(2) per spin
  \addplot[levellinegrey, domain=0:1, samples=2] {0.6931};
  \node[lbl, text=birrengray, font=\scriptsize, anchor=south west]
    at (axis cs:0.025,0.70) {$\log 2$};

  % N -> infinity: the binary entropy
  \addplot[birrenred, line width=1.1pt, domain=0.0008:0.9992, samples=241]
    {-x*ln(x)-(1-x)*ln(1-x)};
  \node[lbl, text=birrenred, anchor=west] at (axis cs:0.70,0.655)
    {$N\to\infty$};

  % N = 24, exact binomials
  \addplot[only marks, mark=*, mark size=0.9pt,
           mark options={draw=birrenorange, fill=birrenorange}]
    coordinates {
      (0.0000,0.0000) (0.0417,0.1324) (0.0833,0.2342) (0.1250,0.3172)
      (0.1667,0.3863) (0.2083,0.4441) (0.2500,0.4921) (0.2917,0.5314)
      (0.3333,0.5628) (0.3750,0.5868) (0.4167,0.6037) (0.4583,0.6138)
      (0.5000,0.6171) (0.5417,0.6138) (0.5833,0.6037) (0.6250,0.5868)
      (0.6667,0.5628) (0.7083,0.5314) (0.7500,0.4921) (0.7917,0.4441)
      (0.8333,0.3863) (0.8750,0.3172) (0.9167,0.2342) (0.9583,0.1324)
      (1.0000,0.0000)
    };
  \node[lbl, text=birrenorange, font=\scriptsize, anchor=north]
    at (axis cs:0.5,0.585) {$N=24$};

  % N = 6 — the system of panel (a).  Filled: the two levels of the shell.
  \addplot[only marks, mark=o, mark size=1.4pt,
           mark options={draw=cgblue, fill=white, line width=0.5pt}]
    coordinates {
      (0.0000,0.0000) (0.1667,0.2986) (0.6667,0.4513) (0.8333,0.2986)
      (1.0000,0.0000)
    };
  \addplot[only marks, mark=*, mark size=1.4pt,
           mark options={draw=cgblue, fill=cgblue}]
    coordinates {(0.3333,0.4513) (0.5000,0.4993)};
  \node[lbl, text=cgblue, font=\scriptsize, anchor=north]
    at (axis cs:0.5,0.425) {$N=6$};

  % What is plotted
  \node[lbl, text=birrengray, font=\scriptsize, anchor=south]
    at (axis cs:0.5,0.10) {\tpcifenglish{shell $=$ one level}%
                                        {Schale $=$ ein Niveau}};
  \node[lbl, text=birrengray, font=\scriptsize, anchor=south]
    at (axis cs:0.5,0.025) {$N_E=\tbinom{N}{k}$, $e=k/N$};
\end{axis}

\end{tikzpicture}
\end{center}
